% TEST LAYOUT (12 November 2026): main.tex with the table split by framework (PyTorch rows, JAX rows,
% hand-written kernel) and the text rewritten around it. Not the submission; main.tex is.
% DMRN+21: the official 2026 kit ships the 2025 style; its event header and
% footer are updated in the local copy, without changing its page layout.
\documentclass[a4paper]{article}
\usepackage[T1]{fontenc}
\usepackage{caption}
\usepackage{times,amsmath,amssymb,graphicx,url,dmrn2026}
\usepackage{indentfirst}
\usepackage[table]{xcolor}
\usepackage{multirow}
\definecolor{heatpurple}{HTML}{DDD0EE}
\captionsetup[figure]{skip=2pt}\setlength{\intextsep}{3pt plus 1pt minus 1pt}
\setlength{\parskip}{0.1em}
\captionsetup[figure]{font=footnotesize,labelfont=footnotesize}
\captionsetup[table]{font=footnotesize,labelfont=footnotesize,skip=2pt}
% TEST LAYOUT: 6pt (main.tex has LaTeX's default 20pt), to make room for the two footnotes at the foot of the right column.
\setlength{\textfloatsep}{6pt plus 2pt minus 4pt}
\setlength{\belowcaptionskip}{0pt}
\renewcommand{\thesection}{\Roman{section}}
\newcommand{\head}[1]{\par\vspace{1pt}\noindent\textbf{#1.}\ \ignorespaces}
% Step arrow: the row or stage builds on the one before it.
\newcommand{\step}{\rotatebox[origin=c]{180}{$\Lsh$}}
% Tighter display spacing. \ninept makes \normalsize an alias of \small, which
% resets the display skips at every size switch, so the reduced values are
% appended to \small itself.
\makeatletter
\g@addto@macro\small{\abovedisplayskip 4pt plus 1pt minus 1pt \belowdisplayskip 4pt plus 1pt minus 1pt
 \abovedisplayshortskip 2pt \belowdisplayshortskip 2pt plus 1pt}
\makeatother
% No extra space between reference entries.
\let\oldthebibliography\thebibliography
\renewcommand{\thebibliography}[1]{\oldthebibliography{#1}\setlength{\itemsep}{0pt}\setlength{\parskip}{0pt}}
% Layout economy after the DMRN+15 auraloss one-pager: a text block 16 mm taller than the
% template's 250 mm. TEST LAYOUT: the title block gaps are 1.2em, 1em, 1em (the template's and main.tex's
% are 2em, 1.5em, 1.5em) to make room for the six extra table rows; the footer distance (0.3in) is the template's.
\setlength{\textheight}{266truemm}\setlength{\footskip}{0.3in}
\cfoot{\fontsize{7}{8}\selectfont\textsc{DMRN+21: DIGITAL MUSIC RESEARCH NETWORK ONE-DAY WORKSHOP 2026, KING'S COLLEGE LONDON, FRI 11 DECEMBER 2026}}
\makeatletter
% A footnote mark set by hand inside another footnote's text, in the style of the automatic ones.
\newcommand{\inlinemark}[1]{{\def\@thefnmark{#1}\normalfont\footnotesize\@makefnmark}}
\def\@maketitle{\newpage
 \null
 \begin{minipage}{0.5\textwidth}
 \textsc{DMRN+21: Digital Music Research Network\\One-day Workshop 2026}\vspace{3mm}\\\textsc{King's College London}\\\textsc{Fri 11 December 2026}
 \end{minipage}%
 \begin{minipage}{0.5\textwidth}
 \hfill\hbox{\includegraphics[width=6cm]{DMRN_Logo_high_res.jpg}}
 \end{minipage}
 \vskip 1.2em \begin{center}
 {\large \bf \@title \par} \vskip 1em {\large
\begin{tabular}[t]{c}\@name \\ \@address
 \end{tabular}
 \par}
 \end{center}
 \par
 \vskip 1em}
\makeatother

\title{Efficient Backpropagation for\\Finite-Difference Nonlinear String Synthesis}
% The acknowledgement is two lines, 8.89pt apart (9.36pt times the 0.95 line spread): the distance between the two
% one-line footnotes at the foot of the right column, so that both footnote rules sit at the same height.
% The DOI may break after its slash: unbroken, it pushes the acknowledgement to three lines.
\name{\normalfont Pablo Tablas de Paula$^{1}$, Rodrigo Diaz$^{1,2}$, Emmanouil Benetos$^{1}$ and Joshua D. Reiss$^{1}$\thanks{\fontsize{7}{9.36}\selectfont Supported by UKRI grant EP/S022694 and QMUL's Apocrita HPC (doi:10.5281/\allowbreak zenodo.438045). We thank Eloi Moliner for advice on text-guided sound design.}}
\address{\small $^{1}$Centre for Digital Music, Queen Mary University of London\quad $^{2}$Fraunhofer HHI, Berlin, Germany\quad p.tablasdepaula@qmul.ac.uk}

\begin{document}
\ninept
\maketitle
\begin{sloppy}
\begin{abstract}
Although realistic, finite-difference time-domain schemes are a rare choice
in differentiable sound synthesis due to their slow and memory-hungry
backpropagation. We compare ways to reduce this cost for a string whose
tension rises as it vibrates, the Kirchhoff--Carrier model. We propose a
method that steps through the string's modes in parallel, with a backward
pass derived by hand for the GPU, and show it is faster and lighter than
automatic differentiation, even when compiled and checkpointed. An online
supplement applies it to sound matching of recorded guitar and
text-guided sound design.\footnotemark
\end{abstract}

\section{Introduction}
Finite-difference time-domain (FDTD) sound synthesis is increasingly lean~\cite{bilbao2024guitar} and runs on the GPU~\cite{lee2024string}, yet differentiable
frameworks~\cite{hayes2024review} largely avoid it because
backpropagating through a time-stepping scheme builds deep recurrent graphs. Hence, the literature has favoured
digital waveguides~\cite{tablas2026sound} and modal
synthesis~\cite{diaz2025}. Physics simulation has long met the same obstacle
with the adjoint method and checkpointing, now the default of
differential-equation solvers~\cite{kidger2021}.
We apply this recipe to a nonlinear FDTD string, propose stepping through
its modes in parallel via a sine transform, and measure what each
ingredient buys in PyTorch\footnotemark{} and in JAX.\footnotemark

\section{Model and method}
\head{Kirchhoff--Carrier string}
The Kirchhoff--Carrier string~\cite[Sec.~8.1]{bilbao2009} is stiff and lossy,
and its tension rises with its stretch. Simply supported at both
ends~\cite{russo2026}, it obeys
\begin{equation}
\begin{aligned}
\rho A\frac{\partial^2u}{\partial t^2}={}&
 \bigg[T_0+\frac{EA}{2L}\!\int_0^L\!\Big(\frac{\partial u}{\partial x}\Big)^{\!2}dx\bigg]
 \frac{\partial^2u}{\partial x^2}
 -EI\frac{\partial^4u}{\partial x^4}\\
 &-2\rho A\Big(\sigma_0-\sigma_1\frac{\partial^2}{\partial x^2}\Big)
 \frac{\partial u}{\partial t},
\end{aligned}
\label{eq:pde}
\end{equation}
with displacement $u(x,t)$, density $\rho$, cross-section $A$, Young's
modulus $E$, second moment of area $I$, length $L$ and damping
$\sigma_0,\sigma_1$. The bracket is the tension: its rest value $T_0$ plus a
term that grows with the string's stretch, one number for the whole string,
which is what bends the pitch of a hard pluck. We integrate
(\ref{eq:pde}) with the energy-conserving scheme of Russo \emph{et
al.}~\cite{russo2026}. A fourth-order five-point
stencil~\cite[Sec.~6.3.1]{bilbao2009} reduces
numerical dispersion.

\head{Backpropagation}
We compare nine implementations in Table~\ref{tab:benchmark}, in PyTorch
and in JAX, each adding one ingredient to the row above. \emph{Auto} is the
scheme's adjoint, derived by automatic differentiation. Every state is stored, and
every operation of every step incurs the overhead of a Python call;
\emph{+~compiled} compiles blocks of 16 unrolled steps in PyTorch, and the
scheme and its gradient as one program in JAX. \emph{+~checkpoints} stores the state only
periodically and re-runs one segment at a time on the reverse pass.
\emph{+~modes} applies a discrete sine transform, after which each stencil
becomes one multiplication per mode and the modes interact only through the
nonlinear tension term. \emph{Hand-written} is the same adjoint with the
same checkpoints and modes, derived by hand and written as a single GPU
(CUDA\footnotemark) kernel, parallel over modes and batch items. All nine return the
gradient of the discrete scheme.

\section{Evaluation}
% The texts of footnotes 1 (abstract), 2 and 3 (Introduction) and 4 (Backpropagation) are placed here so that they sit
% at the foot of the right column. Footnotes 2 to 4 share one line: 3 and 4 are set inside the text of footnote 2,
% with their marks from \inlinemark, so the block stays two lines high. They sit at the foot of the
% right column; the left column then holds the whole Backpropagation paragraph above the acknowledgement.
\footnotetext[1]{\fontsize{6}{7}\selectfont\url{ptablasdpaula.github.io/DMRN21-Kirchhoff_Carrier_String}}%
\footnotetext[2]{\fontsize{6}{7}\selectfont\url{pytorch.org}\quad\inlinemark{3}\url{docs.jax.dev}\quad\inlinemark{4}\url{developer.nvidia.com/cuda}}%
Each evaluation is one sound matching step: a batch of strings is simulated
for 1 s, the mean absolute error of each waveform against the target's is backpropagated to ten parameters:
$\rho$, $T_0$, $r$ ($A=\pi r^2$, $I=\pi r^4/4$), $E$, $L$, $\sigma_0$,
$\sigma_1$, the pluck amplitude and the pluck and pickup positions. The target
is a steel string ($\rho=8050$ kg/m$^3$, $T_0=75$ N, $r=0.356$ mm, $E=174$ GPa,
$L=1$ m, $\sigma_0=0.92$ s$^{-1}$, $\sigma_1=2.9\times10^{-4}$ m$^2$/s), plucked
at its centre with a 2 mm raised cosine of half-width $0.24L$, its
displacement read at its midpoint; the batch perturbs every parameter by up to
$\pm3\%$. The scheme divides the string into 176 equal segments and runs
at 88.2 kHz, decimated to 44.1 kHz, on an A100 80 GB in float64, with a
checkpoint every 256 steps. In Table~\ref{tab:benchmark}, memory
grows linearly with the batch without checkpoints, nearing or exceeding the
A100\textquotesingle s 80\,GB, and stays around 2\,GB or less with them.
Compiling cuts PyTorch\textquotesingle s time by an order of magnitude and
JAX\textquotesingle s by three, making JAX the faster; mode-stepping cuts both further. Our kernel is faster and lighter than the fastest of them.

\begin{table}[t]
 \centering
 \setlength{\tabcolsep}{3.5pt}
 \renewcommand{\arraystretch}{0.76}
 \resizebox{\columnwidth}{!}{% Generated by build_benchmark_summary_table.py from native-rate-two-frameworks-merged.json
\begin{tabular}{@{}c@{\hspace{3pt}}l*{8}{c}}
\multicolumn{2}{@{}l}{} & \multicolumn{8}{c}{Batch Size}\\
\multicolumn{2}{@{}l}{Method} & 1 & 2 & 4 & 8 & 16 & 32 & 64 & 128\\
\hline
\noalign{\vskip 2pt}
\multicolumn{2}{@{}l}{} & \multicolumn{8}{c}{Peak memory (GB)}\\[2pt]
\multirow{4}{*}{\rotatebox[origin=c]{90}{\scriptsize PyTorch}} & Auto & \cellcolor{heatpurple!35.48}2.3 & \cellcolor{heatpurple!41.94}3.8 & \cellcolor{heatpurple!54.84}6.3 & \cellcolor{heatpurple!64.52}11.8 & \cellcolor{heatpurple!77.42}22.8 & \cellcolor{heatpurple!90.32}44.9 & \cellcolor{heatpurple}\textemdash & \cellcolor{heatpurple}\textemdash\\
 & \step\,+ compiled & \cellcolor{heatpurple!25.81}1.4 & \cellcolor{heatpurple!35.48}2.3 & \cellcolor{heatpurple!38.71}3.3 & \cellcolor{heatpurple!48.39}5.4 & \cellcolor{heatpurple!58.06}9.6 & \cellcolor{heatpurple!67.74}18.1 & \cellcolor{heatpurple!80.65}34.9 & \cellcolor{heatpurple!100.00}71.9\\
 & \step\,+ checkpoints & \cellcolor{heatpurple!3.23}0.7 & \cellcolor{heatpurple!12.90}1.0 & \cellcolor{heatpurple!12.90}1.0 & \cellcolor{heatpurple!16.13}1.1 & \cellcolor{heatpurple!16.13}1.1 & \cellcolor{heatpurple!19.35}1.2 & \cellcolor{heatpurple!25.81}1.4 & \cellcolor{heatpurple!29.03}1.8\\
 & \step\,+ modes & \textbf{0.6} & \textbf{0.6} & \textbf{0.6} & \textbf{0.6} & \cellcolor{heatpurple!3.23}\textbf{0.7} & \cellcolor{heatpurple!6.45}0.8 & \cellcolor{heatpurple!9.68}\textbf{0.9} & \cellcolor{heatpurple!22.58}1.3\\
\noalign{\vskip 1.5pt}
\multirow{4}{*}{\rotatebox[origin=c]{90}{\scriptsize JAX}} & Auto & \cellcolor{heatpurple!51.61}6.0 & \cellcolor{heatpurple!51.61}6.0 & \cellcolor{heatpurple!61.29}10.3 & \cellcolor{heatpurple!74.19}18.9 & \cellcolor{heatpurple!87.10}36.1 & \cellcolor{heatpurple!96.77}65.6 & \cellcolor{heatpurple}\textemdash & \cellcolor{heatpurple}\textemdash\\
 & \step\,+ compiled & \cellcolor{heatpurple!45.16}5.0 & \cellcolor{heatpurple!58.06}9.6 & \cellcolor{heatpurple!70.97}18.3 & \cellcolor{heatpurple!83.87}36.0 & \cellcolor{heatpurple!93.55}64.3 & \cellcolor{heatpurple}\textemdash & \cellcolor{heatpurple}\textemdash & \cellcolor{heatpurple}\textemdash\\
 & \step\,+ checkpoints & \cellcolor{heatpurple!3.23}0.7 & \textbf{0.6} & \cellcolor{heatpurple!3.23}0.7 & \textbf{0.6} & \cellcolor{heatpurple!3.23}\textbf{0.7} & \cellcolor{heatpurple!12.90}1.0 & \cellcolor{heatpurple!22.58}1.3 & \cellcolor{heatpurple!32.26}2.1\\
 & \step\,+ modes & \cellcolor{heatpurple!3.23}0.7 & \cellcolor{heatpurple!3.23}0.7 & \cellcolor{heatpurple!3.23}0.7 & \cellcolor{heatpurple!6.45}0.8 & \cellcolor{heatpurple!6.45}0.8 & \cellcolor{heatpurple!12.90}1.0 & \cellcolor{heatpurple!22.58}1.3 & \cellcolor{heatpurple!32.26}2.1\\
\noalign{\vskip 1.5pt}
\multicolumn{2}{@{}l}{Hand-written} & \textbf{0.6} & \textbf{0.6} & \textbf{0.6} & \textbf{0.6} & \cellcolor{heatpurple!3.23}\textbf{0.7} & \cellcolor{heatpurple!3.23}\textbf{0.7} & \cellcolor{heatpurple!9.68}\textbf{0.9} & \cellcolor{heatpurple!19.35}\textbf{1.2}\\
\noalign{\vskip 2pt}\cline{3-10}
\noalign{\vskip 2pt}
\multicolumn{2}{@{}l}{} & \multicolumn{8}{c}{Iteration time (s)}\\[2pt]
\multirow{4}{*}{\rotatebox[origin=c]{90}{\scriptsize PyTorch}} & Auto & \cellcolor{heatpurple!84.62}169 & \cellcolor{heatpurple!84.62}169 & \cellcolor{heatpurple!86.54}171 & \cellcolor{heatpurple!86.54}171 & \cellcolor{heatpurple!88.46}174 & \cellcolor{heatpurple!90.38}177 & \cellcolor{heatpurple}\textemdash & \cellcolor{heatpurple}\textemdash\\
 & \step\,+ compiled & \cellcolor{heatpurple!57.69}10.7 & \cellcolor{heatpurple!59.62}11.0 & \cellcolor{heatpurple!61.54}11.2 & \cellcolor{heatpurple!63.46}11.3 & \cellcolor{heatpurple!65.38}11.8 & \cellcolor{heatpurple!65.38}11.8 & \cellcolor{heatpurple!67.31}12.1 & \cellcolor{heatpurple!69.23}12.4\\
 & \step\,+ checkpoints & \cellcolor{heatpurple!71.15}13.2 & \cellcolor{heatpurple!75.00}13.6 & \cellcolor{heatpurple!80.77}13.9 & \cellcolor{heatpurple!76.92}13.7 & \cellcolor{heatpurple!73.08}13.5 & \cellcolor{heatpurple!80.77}13.9 & \cellcolor{heatpurple!78.85}13.8 & \cellcolor{heatpurple!82.69}14.2\\
 & \step\,+ modes & \cellcolor{heatpurple!40.38}4.31 & \cellcolor{heatpurple!44.23}4.52 & \cellcolor{heatpurple!44.23}4.52 & \cellcolor{heatpurple!51.92}4.59 & \cellcolor{heatpurple!46.15}4.54 & \cellcolor{heatpurple!53.85}4.60 & \cellcolor{heatpurple!48.08}4.55 & \cellcolor{heatpurple!55.77}4.68\\
\noalign{\vskip 1.5pt}
\multirow{4}{*}{\rotatebox[origin=c]{90}{\scriptsize JAX}} & Auto & \cellcolor{heatpurple!94.23}5324 & \cellcolor{heatpurple!96.15}5355 & \cellcolor{heatpurple!98.08}5364 & \cellcolor{heatpurple!92.31}5315 & \cellcolor{heatpurple!100.00}5374 & \cellcolor{heatpurple!92.31}5315 & \cellcolor{heatpurple}\textemdash & \cellcolor{heatpurple}\textemdash\\
 & \step\,+ compiled & \cellcolor{heatpurple!13.46}2.41 & \cellcolor{heatpurple!9.62}2.26 & \cellcolor{heatpurple!15.38}2.49 & \cellcolor{heatpurple!19.23}2.99 & \cellcolor{heatpurple!23.08}3.13 & \cellcolor{heatpurple}\textemdash & \cellcolor{heatpurple}\textemdash & \cellcolor{heatpurple}\textemdash\\
 & \step\,+ checkpoints & \cellcolor{heatpurple!21.15}3.04 & \cellcolor{heatpurple!17.31}2.68 & \cellcolor{heatpurple!25.00}3.18 & \cellcolor{heatpurple!34.62}4.01 & \cellcolor{heatpurple!38.46}4.25 & \cellcolor{heatpurple!36.54}4.23 & \cellcolor{heatpurple!42.31}4.45 & \cellcolor{heatpurple!50.00}4.56\\
 & \step\,+ modes & \cellcolor{heatpurple!7.69}2.02 & \cellcolor{heatpurple!5.77}1.88 & \cellcolor{heatpurple!11.54}2.34 & \cellcolor{heatpurple!21.15}3.04 & \cellcolor{heatpurple!26.92}3.22 & \cellcolor{heatpurple!28.85}3.33 & \cellcolor{heatpurple!30.77}3.51 & \cellcolor{heatpurple!32.69}3.65\\
\noalign{\vskip 1.5pt}
\multicolumn{2}{@{}l}{Hand-written} & \textbf{0.37} & \textbf{0.37} & \textbf{0.37} & \textbf{0.37} & \cellcolor{heatpurple!1.92}\textbf{0.38} & \cellcolor{heatpurple!1.92}\textbf{0.38} & \cellcolor{heatpurple!1.92}\textbf{0.38} & \cellcolor{heatpurple!3.85}\textbf{0.51}\\
\hline
\end{tabular}
}
 \caption{Peak memory and time of one gradient step. \step: adds one ingredient to the row above; \textemdash: exceeds the A100\textquotesingle s 80\,GB.}
 \label{tab:benchmark}
\end{table}

{\footnotesize
\begin{thebibliography}{10}
\bibitem{bilbao2024guitar} S.~Bilbao \emph{et~al.}, ``Real-time guitar synthesis,'' \emph{Proc. DAFx}, 2024.
\bibitem{lee2024string} J.~W.~Lee \emph{et~al.}, ``String sound synthesizer on GPU-accelerated finite difference scheme,'' \emph{Proc. ICASSP}, 2024.
\bibitem{hayes2024review} B.~Hayes \emph{et~al.}, ``A review of differentiable digital signal processing for music and speech synthesis,'' \emph{Front. Sig. Proc.}, v.~3, 1284100, 2024.
\bibitem{tablas2026sound} P.~Tablas~de~Paula \emph{et~al.}, ``Sound matching with a differentiable Karplus-Strong algorithm,'' \emph{Proc. DAFx}, 2026.
\bibitem{diaz2025} R.~Diaz and M.~Sandler, ``Fast differentiable modal simulation of non-linear strings, membranes, and plates,'' \emph{Proc. DAFx}, 2025.
\bibitem{kidger2021} P.~Kidger, ``On neural differential equations,'' DPhil thesis, Oxford, 2021.
\bibitem{bilbao2009} S.~Bilbao, \emph{Numerical Sound Synthesis: Finite Difference Schemes and Simulation in Musical Acoustics}, Wiley, 2009.
\bibitem{russo2026} R.~Russo \emph{et~al.}, ``Numerical convergence of the scalar auxiliary variable method applied to nonlinear stiff string models,'' \emph{Nonlinear Dyn.}, v.~114, 2026.
\end{thebibliography}}
\end{sloppy}
\end{document}
